
\فصل{نمای کلی محصول}

\قسمت{کاربران و نقش‌ها در سامانه}

جدول~\ref{جدول:نقش-دسترسی} ارتباط هر نوع کاربر با سامانه‌ی ساکنا را نشان می‌دهد.

\begin{table}[ht]
\centering
\caption{کاربران ساکنا و دسترسی‌های آن‌ها}
\label{جدول:نقش-دسترسی}
\begin{tabular}{|>{\raggedleft\arraybackslash}p{3cm}|>{\raggedleft\arraybackslash}p{9cm}|}
\hline
\rl{\textbf{کاربر}} & \rl{\textbf{دسترسی‌های اصلی در ساکنا}} \\
\hline
\rl{ساکن} & \rl{مشاهده‌ی اطلاعات واحد، ثبت و پیگیری درخواست تعمیرات، پرداخت شارژ از طریق کیف پول، رزرو و لغو رزرو امکانات، مشاهده‌ی اطلاعیه‌ها} \\
\hline
\rl{مدیر ساختمان} & \rl{مدیریت ساختمان و واحدها، بررسی و ارجاع درخواست‌های خدماتی، تعریف شارژ و هزینه‌ها، مدیریت امکانات مشترک، انتشار اطلاعیه، گزارش‌گیری کامل} \\
\hline
\rl{کارکن خدماتی} & \rl{مشاهده‌ی وظایف ارجاع‌شده، ثبت گزارش انجام کار، مشاهده‌ی تاریخچه‌ی فعالیت‌ها} \\
\hline
\rl{مدیر سامانه} & \rl{مشاهده‌ی فهرست کاربران، غیرفعال‌سازی حساب‌های مشکل‌دار} \\
\hline
\end{tabular}
\end{table}

\قسمت{نمای کلی ارتباط کاربران با سیستم}

\begin{table}[ht]
\centering
\caption{نمای کلی سامانه‌ی ساکنا}
\label{جدول:نمای-کلی}
\renewcommand{\arraystretch}{1.8}
\begin{tabular}{|>{\centering\arraybackslash}p{2.5cm}|>{\centering\arraybackslash}p{7cm}|>{\centering\arraybackslash}p{2.5cm}|}
\hline
 & \rl{\Large\textbf{ساکنا}} & \\
\rl{\textbf{ساکن}} & \rl{\textit{سامانه‌ی آنلاین مدیریت مجتمع مسکونی}} & \rl{\textbf{مدیر ساختمان}} \\
\hline
\rl{احراز هویت \& کیف پول} & \rl{درخواست خدماتی \quad | \quad شارژ و پرداخت \quad | \quad رزرو امکانات} & \rl{مدیریت و گزارش} \\
\hline
\rl{\textbf{کارکن خدماتی}} & \rl{اطلاعیه \quad | \quad احراز هویت \quad | \quad \lr{RBAC}} & \rl{\textbf{مدیر سامانه}} \\
\hline
\end{tabular}
\renewcommand{\arraystretch}{1}
\end{table}

\قسمت{ماژول‌های اصلی سامانه}

ساکنا از ماژول‌های زیر تشکیل شده است:

\begin{table}[ht]
\centering
\caption{ماژول‌های اصلی ساکنا}
\label{جدول:ماژول}
\begin{tabular}{|>{\raggedleft\arraybackslash}p{3.5cm}|>{\raggedleft\arraybackslash}p{3cm}|>{\raggedleft\arraybackslash}p{5cm}|}
\hline
\rl{\textbf{ماژول}} & \rl{\textbf{کاربران}} & \rl{\textbf{قابلیت‌های اصلی}} \\
\hline
\rl{احراز هویت و دسترسی} & \rl{همه} & \rl{ثبت‌نام، ورود، \lr{RBAC}، مدیریت پروفایل} \\
\hline
\rl{مدیریت ساختمان} & \rl{مدیر} & \rl{واحدها، ساکنین، کارکنان} \\
\hline
\rl{درخواست خدماتی} & \rl{ساکن، مدیر، کارکن} & \rl{ثبت، ارجاع، پیگیری، گزارش} \\
\hline
\rl{شارژ و کیف پول} & \rl{ساکن، مدیر} & \rl{تعریف شارژ، پرداخت، سوابق} \\
\hline
\rl{رزرو امکانات} & \rl{ساکن، مدیر} & \rl{تعریف امکان، رزرو، لغو} \\
\hline
\rl{اطلاعیه} & \rl{ساکن، مدیر} & \rl{انتشار، مشاهده، آرشیو} \\
\hline
\rl{گزارش‌گیری} & \rl{مدیر} & \rl{پرداخت‌ها، درخواست‌ها، رزروها} \\
\hline
\end{tabular}
\end{table}

\قسمت{جریان کاری نمونه}

\زیرقسمت{جریان پرداخت شارژ}

\شروع{شمارش}
\فقره مدیر دوره‌ی شارژ را تعریف می‌کند.
\فقره ساکنین اطلاعیه‌ی شارژ را مشاهده می‌کنند.
\فقره ساکن از طریق کیف پول خود مبلغ را پرداخت می‌کند.
\فقره سامانه تأیید پرداخت را ثبت می‌کند.
\فقره مدیر در داشبورد وضعیت پرداخت‌ها را مشاهده می‌کند.
\پایان{شمارش}

\زیرقسمت{جریان درخواست تعمیرات}

\شروع{شمارش}
\فقره ساکن درخواست تعمیرات را با توضیحات ثبت می‌کند.
\فقره مدیر درخواست را بررسی و تأیید می‌کند.
\فقره مدیر درخواست را به کارکن خدماتی مناسب ارجاع می‌دهد.
\فقره کارکن درخواست را مشاهده کرده و کار را انجام می‌دهد.
\فقره کارکن گزارش انجام کار را ثبت می‌کند.
\فقره ساکن از تکمیل درخواست مطلع می‌شود.
\پایان{شمارش}
